% Part: history
% Chapter: biographies 
% Section: alan-turing 

\documentclass[../../../include/open-logic-section]{subfiles}

\begin{document}

\olfileid{his}{bio}{tur}

\olsection{艾伦·Turing}

艾伦·图灵于1912年6月23日出生于伦敦的梅达维尔。他被誉为理论计算机科学之父。图灵自幼便对物理科学和数学产生了兴趣。然而，当时的学校更重视文学和古典学，他儿时的兴趣在学校里并未得到重视。因此，他在学校表现不佳，屡遭老师训斥。

\olphoto{turing-alan}{Alan Turing}

图灵在剑桥大学国王学院攻读本科，学习数学。1936年，图灵提出了（如今所谓的）图灵机，试图精确定义可计算函数的概念，并证明判定问题的不可判定性。但 邱奇通过其提出的 λ 演算率先证明了这一结果。图灵的论文在发表时仍引用了 邱奇 的结果。邱奇 邀请图灵前往普林斯顿，图灵于1936至1938年间在此度过，并在 邱奇 的指导下获得了博士学位。

尽管已投身逻辑学，图灵早年对物理科学的兴趣依然不减。第二次世界大战期间，他在布莱切利园的英国密码分析部门服役，他的实践技能在此得到了发挥。图灵是破解德国海军通信密码（恩尼格玛密码）的核心人物。图灵在统计学和密码学方面的专长，加之电子设备的引入，使得团队能够制造出一台名为“炸弹”的解密机，从而破解了该密码。他的想法也促成了世界上第一台可编程电子计算机“巨人”的诞生，这台计算机同样在布莱切利园用于破解德国的洛伦兹密码。

图灵是同性恋者。尽管如此，1942年他仍向在布莱切利园的一位同事 Joan Clarke 求婚，但后来解除了婚约，并向她坦白了自己的性取向。他一生中有过几段恋情，尽管当时的英国将同性恋行为定性为刑事犯罪。1952年，图灵的住宅被他当时恋人的一个朋友入室盗窃；在报案时，图灵承认了同性关系，因为他误以为政府即将把同性恋行为合法化。事实并非如此，他被指控犯有严重猥亵罪。图灵没有选择入狱，而是选择了接受会导致性欲减退的激素治疗。1954年6月8日，图灵被发现死于氰化物中毒——很可能是自杀。他于2013年获得了伊丽莎白二世女王的皇家赦免。

\begin{reading}
关于艾伦·图灵的全面传记，参见 \citet{Hodges2014}。
图灵的生平和工作启发了一部戏剧《破译密码》，该剧于1996年制作成电视电影，由德里克·雅各比饰演图灵。获得奥斯卡奖提名的电影《模仿游戏》（本尼迪克特·康伯巴奇和凯拉·奈特莉主演）也大致基于艾伦·图灵在布莱切利园的生平和时光 \citep{Imitation2014}。

\citet{Radiolab2012} 有几期关于图灵生平和工作的播客。BBC《地平线》系列纪录片《图灵博士奇异的一生与死亡》可以在线观看 \citep{Sykes1992}。
\citep{Theelen2012} 是一段乐高图灵机运行的视频短片——为纪念图灵2012年的百年诞辰而制作。

图灵关于图灵机与判定问题的最初论文是 \citet{Turing1937}。
\end{reading}

\end{document}